%=====================================================================
%  Part 3 opener --- Research Designs for Computing
%=====================================================================
\structuralfigures{P3.}

This is the catalogue. Each design answers some questions and is
incapable of answering others, and the skill this part teaches is not knowing
all eight --- it is choosing correctly under constraint, and being able to say
what your choice cannot establish. Read it with your own question from
Chapter~3 in front of you. Every chapter ends by telling you what it forecloses.

\begin{conceptbox}[What Part~3 Covers]
\begin{tight}
  \item \textbf{Chapter 9 --- Design Science Research.} Artefacts as contributions; Hevner's guidelines, the DSRM process, and FEDS evaluation.
  \item \textbf{Chapter 10 --- Controlled Experiments.} Randomisation, factorial and within-subjects designs, blocking, and human subjects in computing.
  \item \textbf{Chapter 11 --- Quasi-Experiments and Repository Mining.} Identification strategies when you cannot randomise, and the confounders in observational software data.
  \item \textbf{Chapter 12 --- Case Study Research.} Units of analysis, propositions, chain of evidence, and why a case study is not a small survey.
  \item \textbf{Chapter 13 --- Survey Research.} Sampling frames, instrument construction, piloting, and the nonresponse problem.
  \item \textbf{Chapter 14 --- Qualitative Methods and Coding.} Interviews, grounded theory, thematic analysis, coding reliability and saturation.
  \item \textbf{Chapter 15 --- Mixed Methods.} Convergent and sequential designs, joint displays, and integration that is more than juxtaposition.
  \item \textbf{Chapter 16 --- Benchmarking and Performance Evaluation.} Workload characterisation, statistically defensible measurement, and simulation validity.
\end{tight}
\end{conceptbox}
